\documentclass[11pt]{article}
\usepackage[T1]{fontenc}
\usepackage{lmodern,amsmath,amssymb,amsthm,booktabs,microtype}
\usepackage[margin=28mm]{geometry}
\usepackage[hidelinks]{hyperref}
\hypersetup{pdftitle={Perron defects under reciprocal edge twists},pdfauthor={1200}}
\newtheorem{theorem}{Theorem}
\DeclareMathOperator{\arsinh}{arsinh}
\DeclareMathOperator{\arcosh}{arcosh}
\DeclareMathOperator{\diag}{diag}
\newcommand{\R}{\mathbb R}
\newcommand{\ip}[2]{\langle #1,#2\rangle}
\title{Perron defects under reciprocal edge twists\\
\large Cycle rigidity and minimum realization sizes}
\author{1200}
\date{5 November 2026}
\begin{document}
\maketitle
\begin{abstract}
We compare the Perron root of a reciprocal edge twist of a symmetric
nonnegative matrix with an optimized lower bound at the untwisted Perron
vector. A balanced gauge makes the left and right residuals equal, giving a
Schur-complement formula for their spectral defect. Every finite nonzero
order of vanishing is a multiple of four. On unicyclic support the only
nonzero orders are \(4,8,12\); the zero-defect class at nonzero circulation
is exactly a pure cycle with constant diagonal and two-periodic edge
weights. Every order \(4m\) is realizable in dimension \(2m+2\), which
is sharp in a specified two-cell class. The first three unrestricted
minimum sizes are \(3,3,4\); for zero-diagonal bipartite realizations the
first two are \(4,5\).
\end{abstract}

\section{Setup and relation to prior work}
Let \(B=B^\top\geq0\) entrywise be \(n\times n\), \(n\geq2\), with connected
simple off-diagonal support \(G=(V,E)\). Put \(\lambda_0=\rho(B)\) and
normalize the positive Perron vector \(u\) by \(u^\top u=1\). For
\(\{i,j\}\in E\) let \(c_{ij}=B_{ij}u_i u_j\). A real field \(h\) on
oriented edges is antisymmetric. The twist \(A(h)\) keeps the diagonal
of \(B\) and has off-diagonal entries \(B_{ij}e^{h_{ij}}\) on the support.
Use \((\nabla\phi)_{ij}=\phi_j-\phi_i\) and set
\begin{equation}
 F_h(\phi)=2\sum_{e\in E}c_e\{\cosh(h_e+\nabla_e\phi)-1\},\qquad
 E(h)=\min_{\phi_1=0} F_h(\phi).
 \label{eq:energy}
\end{equation}
Each edge sum counts an undirected edge once; the choice of orientation
does not affect it. Neighbor sums are directed. Define
\(\Delta(t)=\rho(A(th))-\lambda_0-E(th)\).

The starting variational bound is a fixed-vector specialization of the
spectral-radius formula of Altenberg~\cite{altenberg}. The divergence-free
\(\sinh\) balance is related to the force/current decomposition of
Kaiser, Jack and Zimmer~\cite{kaiser}. In Markov coordinates
\(\diag(u)^{-1}B\diag(u)-\lambda_0I\), the stationary law is \(u_i^2\);
fixed-density flow optimization and comparison with driven rates fit the
variational framework of Chetrite and Touchette~\cite[Appendix B]{chetrite}.
Fourth-order perturbation uses the same reduced-resolvent machinery as
current-cumulant expansions~\cite{baiesi}. Related current-fluctuation bounds
and localization of complex generator eigenvalues
\cite{pietzonka,kolchinsky} concern different comparisons. Our focus is the
order, rigidity and realization size of the particular deficit \(\Delta\).
We make no claim of a new general variational principle or of global priority.

\section{Minimizing gauge and perturbation}
For a minimizing field \(z=h+\nabla\phi\), write
\[
 k_i=u_i^{-2}\sum_j c_{ij}(\cosh z_{ij}-1).
\]
\begin{theorem}[T1]
The pinned minimizer in \eqref{eq:energy} exists and is unique. It satisfies
\(\sum_j c_{ij}\sinh z_{ij}=0\) at every vertex.
One has \(\rho(A(h))\geq\lambda_0+E(h)\), with equality exactly when all
\(k_i\) coincide.
\end{theorem}
\begin{proof}
The Hessian in direction \(v\) is
\(2\sum_E c_e\cosh(z_e)(\nabla_e v)^2\). Connectedness makes it positive
on nonzero pinned directions. On a sublevel \(F_h\leq K\), every edge
satisfies \(|h_e+\nabla_e\phi|\leq\arcosh(1+K/(2c_e))\). Summing the
resulting bounds for \(|\nabla_e\phi|\) along fixed paths from vertex 1
bounds every \(\phi_i\). Thus the nonempty sublevel \(F_h\leq F_h(0)\)
is compact, and strict convexity gives uniqueness. The pinned gradient
equations give the divergence condition at all but the pinned vertex;
antisymmetry gives it there too.

Gauge changes are diagonal similarities:
\(A(h+\nabla\phi)=D^{-1}A(h)D\), \(D=\diag(e^{\phi_i})\).
Write a positive right Perron vector of \(A(h)\) as \(u_i e^{q_i}\),
scaled so that \(q_1=0\). Divide its equation in row \(i\) by that
entry, multiply by \(u_i^2\), and sum. Pairing the two orientations
gives \(\rho(A(h))-\lambda_0=F_h(q)\).
At the minimizing gauge, cancellation of the \(\sinh\) terms gives
\((A(z)u)_i/u_i=\lambda_0+k_i\), and \(\sum_i u_i^2 k_i=E(h)\).
If the increments coincide, positivity identifies the displayed
eigenvalue as the Perron root. Conversely equality makes \(q\) the
minimizing gauge, so the Perron equation forces the increments to coincide.
\end{proof}

Let \(\alpha=h+\nabla\psi\) be the harmonic field, characterized by
\(\sum_j c_{ij}\alpha_{ij}=0\) with \(\psi_1=0\); it is the unique
minimizer of the positive grounded quadratic energy.
For any potential \(f\), \(\sum_E c_e\alpha_e\nabla_e f=0\).
Let \(R\) be the inverse of \(\lambda_0 I-B\) on \(u^\perp\), extended
by zero on \(u\). Symmetry and simplicity of the largest eigenvalue give
strict positivity of \(R\) on \(u^\perp\), including in the bipartite case.

\begin{theorem}[T2]
Put \(K_{2,ij}=B_{ij}\alpha_{ij}^2/2\) off the diagonal and \(K_{2,ii}=0\).
For \(a=\sum_E c_e\alpha_e^2\), \(b=\sum_E c_e\alpha_e^4/12\), and
\(g=K_2u-au\),
\begin{align*}
 E(th)&=at^2+bt^4+O(t^6),\\
 \rho(A(th))&=\lambda_0+at^2+(b+g^\top Rg)t^4+O(t^6).
\end{align*}
The fourth-order coefficient of \(\Delta\) is nonnegative, vanishing
exactly when \(K_2u=au\). The quantities are invariant under changing \(h\)
by a gradient.
\end{theorem}
\begin{proof}
Use the harmonic gauge. The divergence equations for the pinned
minimizer \(\eta(t)\) at field \(t\alpha\) are analytic and their
linearization in \(\eta\) at zero is the negative grounded weighted
Laplacian. The analytic implicit-function theorem and uniqueness
give an analytic odd solution, with \(\eta'(0)=0\) by harmonicity.
In the energy expansion the fourth-order mixed term
\(2\sum_E c_e\alpha_e\nabla_e\eta_3\) vanishes. This yields \(a,b\).

Write \(A(t\alpha)=B+\sum_{\ell\geq1}t^\ell K_\ell\), where
\(K_{\ell,ij}=B_{ij}\alpha_{ij}^{\ell}/\ell!\) off the diagonal.
Odd \(K_\ell\) are skew, even ones symmetric; \(K_1u=0=u^\top K_1\).
The simple eigenpair is analytic near zero, with eigenvector normalized
by \(u^\top r(t)=1\), positive near zero. Transposition makes the
eigenvalue even. Comparing orders one and two gives \(r_1=0\),
\(u^\top K_2u=a\), and \(r_2=Rg\). Multiplying the fourth-order
equation by \(u^\top\) eliminates \(K_1r_3\) and all eigenvector
normalization terms. It leaves \(u^\top K_4u+u^\top K_2Rg=b+g^\top Rg\).
Finally \(u^\top g=0\) and \(R\) is positive on \(u^\perp\).
\end{proof}

\section{Cycle current and the order mechanism}
\begin{theorem}[T3]
On a unicyclic support with oriented cycle \(C\) and circulation
\(\Phi=\sum_C h_e\), the unique \(J\in\R\) satisfying
\(\sum_C\arsinh(J/c_e)=\Phi\) gives
\[
 E(h)=2\sum_{e\in C}(\sqrt{c_e^2+J^2}-c_e).
\]
At nonzero \(\Phi\), any vertex off the cycle forces a strict bound in T1.
At \(\Phi=0\), the twist is diagonally similar to \(B\).
\end{theorem}
\begin{proof}
The divergence-free current \(c_{ij}\sinh z_{ij}\) vanishes on a bridge
by summing divergences on one side. It is therefore constant in the cycle
orientation. The sum of \(\arsinh(J/c_e)\) is strictly increasing and
onto. This choice of field has the right cycle sum and differs from \(h\)
by a gradient, as seen by integrating on a spanning tree. Substitution
gives the energy. At nonzero circulation the increments are positive
on the cycle and zero outside, so they cannot all coincide. Zero
circulation gives \(J=0\).
\end{proof}

\begin{theorem}[T4]
For the minimizing \(z(t)\) at \(th\) set \(e(t)=E(th)\) and
\(v_i(t)=u_i(k_i(t)-e(t))\). The functions \(e,k,v\) are even and
real analytic near zero. If \(v\equiv0\) near zero then \(\Delta\equiv0\)
there. Otherwise let \(2m\) be the least degree of a nonzero Taylor
coefficient \(v_{2m}\). Then
\begin{equation}
 \Delta(t)=(v_{2m}^{\top}R v_{2m})\,t^{4m}+O(t^{4m+2}),
 \qquad v_{2m}^{\top}R v_{2m}>0 .
 \label{eq:hierarchy}
\end{equation}
\end{theorem}
\begin{proof}
The minimizing field is analytic and odd by the implicit-function
argument and gauge invariance. Put \(M(t)=A(z(t))\). Balance cancels
\(\sinh z\) in both \(M(t)u\) and \(M(t)^\top u\). Thus relative to
\(\R u\oplus u^\perp\) the matrix has block form
\[
 M(t)=\begin{pmatrix}\lambda_0+e(t)&v(t)^\top\\ v(t)&\mathcal C(t)\end{pmatrix}.
\]
In particular \(v\perp u\). The compression \(\mathcal C(0)\) has no
eigenvalue \(\lambda_0\), so \(\lambda(t)I-\mathcal C(t)\) is invertible
near zero for \(\lambda(t)=\rho(A(th))\). A positive right eigenvector
has nonzero \(u\)-component; eliminating its perpendicular component yields
\[
 \Delta(t)=v(t)^\top(\lambda(t)I-\mathcal C(t))^{-1}v(t).
\]
At zero the inverse is \(R|_{u^\perp}\). Expanding gives
\eqref{eq:hierarchy}, with the odd remainder coefficient absent by
transposition symmetry. In particular \(v_2=g\); if it vanishes
then \(\Delta=O(t^8)\).
\end{proof}

\section{Sharp rigidity on one cycle}
\begin{theorem}[T5]
For unicyclic support and nonzero circulation, the defect has leading
order \(4,8,12\), with positive coefficient, or vanishes for every real
\(t\). The last case occurs exactly for a pure cycle with a constant
diagonal and two-periodic edge weights. Odd cycles then have uniform
weights. Any vertex outside the cycle forces order four, and all three
finite orders occur.
\end{theorem}
\begin{proof}
At an outside vertex the harmonic field vanishes on incident bridges,
so \(g_i=-au_i\ne0\). On a pure cycle write \(r_i=1/c_i\) for successive
edge resistances. The current is an analytic local parameter since
\(J'(0)=\Phi/\sum_i r_i\ne0\). The first three coefficients of
\(k_i(J)\) at degrees \(2,4,6\) are the incident sums of \(r,r^3,r^5\)
divided by \(u_i^2\), multiplied by \(1/2,-1/8,1/16\), respectively.
Changing from \(J\) to \(t\) preserves the first degree with a
nonconstant coefficient.

If all three coefficients are constant across vertices, then \(u_i^2\)
is proportional to \(r_{i-1}+r_i\). Ratios of the higher incident sums
to the first yield common values on every adjacent pair \((x,y)\):
\[
 Q=x^2-xy+y^2,\qquad H=(x^5+y^5)/(x+y)=Q^2+xyQ-(xy)^2.
\]
Compare \((x,y)\) and \((y,z)\). Equality of \(Q\) implies
\((x-z)(x+z-y)=0\). If \(x\ne z\) then \(y=x+z\); the difference of
the \(H\) values is \(y(x-z)(Q-y^2)=-xyz(x-z)\ne0\), impossible.
Thus \(r_{i+1}=r_{i-1}\), and the first test makes \(u\) uniform.
The Perron equation gives constant diagonal and two-periodic weights
of \(B\). Conversely those weights and a constant diagonal yield
uniform \(u\) and identical increments at every \(J\); T1 then gives
global equality. If any of the three tests fails, the first such test
and T4 give the stated order. The following exact witnesses prove sharpness.
\end{proof}

\paragraph{Exact witnesses.}
Use the unnormalized convention \(U>0\), \(M=U^\top U\), and positive
conductances \(C_i\) around a pure cycle. Define
\(B_{ij}=C_{ij}/(U_iU_j)\) on edges and
\(B_{ii}=\lambda_0-\sum_j C_{ij}/U_i^2\).
In each row below the diagonal is nonnegative and \(BU=\lambda_0U\).
Write \(r_i=1/C_i\), \(S=\sum_i r_i\); set \(h_i=r_i/S\) in the
cycle orientation. The conductances in \eqref{eq:energy} are \(C_i/M\).
\begin{center}
\small
\begin{tabular}{@{}cllc@{}}
\toprule
Order & \(U\) & \(C\) around the cycle & \(\lambda_0\)\\
\midrule
\(4\)  & \((1,1,1)\) & \((1,1,2)\) & \(5\)\\
\(8\)  & \((2,3,3)\) & \((1/2,1/7,1/2)\) & \(1\)\\
\(12\) & \((3,3,4,4)\) & \((3/2,3/25,3/23,3/25)\) & \(1\)\\
\bottomrule
\end{tabular}
\end{center}
For the first nonconstant incident moment
\(F_i=(r_{i-1}^{2j-1}+r_i^{2j-1})/U_i^2\), the leading
\(\sqrt M\,v_{2j}\) is \(b_j U_i(F_i-\bar F)/S^{2j}\) in row \(i\),
where \(\bar F=\sum_i U_i^2 F_i/M\) and
\(b_1=1/2,b_2=-1/8,b_3=1/16\).
Indeed the lower moments are constant and cancel after centering,
and the current for \(C\) has first coefficient \(1/S\).
In the eighth-order example \(S=11\), the first moment is constant and
the third is \((4,39,39)\). In the twelfth-order example \(S=25\),
the first and third are constant, \(1\) and \(193/3\); the fifth
takes two values whose difference is \(8050/27\).
Substitution gives these certificates (each \(x\perp U\)):
\[
\begin{array}{c|c|c|c}
 4j & x & \sqrt M\,v_{2j} & Rx\\ \hline
 4 & (-1,2,-1) & x/75 & x/3\\
 8 & (3,-1,-1) & 105x/(4\cdot11^5) & 36x/11\\
 12 & (4,4,-3,-3) & 161x/(18\cdot25^6) & 48x
\end{array}
\]
The coefficients from \eqref{eq:hierarchy} are respectively
\(2/16875,\ 99225/(8\cdot11^{11}),\ 103684/(27\cdot25^{12})\).
The resolvent relations are verified by multiplying by
\(\lambda_0 I-B\); the factor \(1/M\) in the quadratic form converts
the unnormalized vectors to unit-Perron coordinates.

\section{Every multiple of four and smaller realizations}
\begin{theorem}[T6--T7]
For every integer \(m\geq1\), rational matrices with positive diagonal
and integer twists realize the coefficient \(1/[2((2m)!)^2]\) at order
\(4m\), in dimensions \(4m+2\) and \(2m+2\). Each construction has a
connected zero-diagonal bipartite lift preserving the entire defect,
in dimensions \(8m+4\) and \(4m+4\), respectively.
The dimension \(2m+2\) is minimal within the two-cell class specified below.
\end{theorem}
\begin{proof}
For \(1\leq a\leq m\) put
\[
 d_a=\frac{1}{a^2\prod_{b\ne a}(a^2-b^2)},\qquad
 \sum_{a=1}^m d_a a^{2j}=\begin{cases}0,&1\leq j<m,\\1,&j=m.\end{cases}
\]
The identity is the highest-coefficient comparison for Lagrange
interpolation of \(x^{j-1}\) at the distinct nodes \(a^2\).
Let \(D(t)=\sum_a d_a(\cosh(at)-1)=t^{2m}/(2m)!+O(t^{2m+2})\).

For the first construction take two cyclic cells of size \(2m+1\).
The oriented edges \(i\to i+a\) for all vertices \(i\) and \(1\leq a\leq m\)
cover each unordered pair exactly once. Use twist \(a\) and weights
\(w_a^\pm=1+\max(\pm d_a,0)\) in the two cells. Connect corresponding
vertices by weight-one edges of zero twist. Choose a common row sum
\(\lambda_0=2+2\sum_a\max(w_a^+,w_a^-)\) by filling the diagonal.
All diagonal entries are at least one. The Perron vector is uniform.
At zero potential, each color balances in the \(\sinh\) equation at
every vertex. The row increments \(k_\pm=2\sum_a w_a^\pm(\cosh(at)-1)\)
are constant on the cells and \(E=(k_++k_-)/2\).
The cell-constant quotient has diagonal
\(\lambda_0-1+k_\pm\) and off-diagonal entries one. Its positive
Perron eigenvector lifts, so its Perron root is the full one. Therefore
\(\Delta(t)=\sqrt{1+D(t)^2}-1\), proving the claimed order and constant.

For the smaller construction put the positive and negative indices of
\(d_a\) in separate cells, with \(k_+,k_-\) indices respectively and
\(L_\pm=2k_\pm+1\) vertices in each cell. The \(r\)-th index \(a\) in a
cell is assigned to all cyclic edges \(i\to i+r\), with weight
\(|d_a|\) and twist \(a\). A cell without indices is a singleton.
Write \(n=L_++L_-\), \(p=L_+/n\), \(q=L_-/n\), and give every
cross-cell edge weight \(w=8L_+L_-/n^3\) and twist zero.
Then \(\gamma=nw=8pq\), with cross-cell row sums \(q\gamma,p\gamma\).
Choose \(\lambda_0\) one larger than the maximum off-diagonal row
sum and fill the diagonal. Again the Perron vector is uniform, zero
potential minimizes the energy, and \(E=p k_++q k_-\).
Now \(d=k_+-k_-=2D\). The quotient after subtracting
\((\lambda_0+E)I\) is
\[
 \begin{pmatrix} -q\gamma+qd & q\gamma\\ p\gamma &-p\gamma-pd\end{pmatrix}.
\]
Its trace is \(-\gamma+(q-p)d\) and determinant \(-pq d^2\).
The nonnegative root \(\Delta\) satisfies
\(\Delta^2+[\gamma+(p-q)d]\Delta-pq d^2=0\).
Expansion of the larger root near zero gives
\(\Delta=(pq/\gamma)d^2+O(d^3)\), and the moment identity yields
the same coefficient \(1/[2((2m)!)^2]\) and remainder \(O(t^{4m+2})\).

For either construction use the bipartite lift
\(\widehat B=\left(\begin{smallmatrix}0&B\\ B&0\end{smallmatrix}\right)\)
with off-diagonal blocks \(A(th)\) in the twist. The edges from the
positive diagonal connect corresponding vertices across sheets; the
original support ensures connectedness. A positive Perron vector
\(r(t)\) lifts to \((r(t),r(t))\), so the root is unchanged.
The normalized base vector is \((u,u)/\sqrt2\). The energy and
its minimizer modulo constants are invariant under swapping sheets;
if the potentials differed by a constant under this involution, two
swaps would double it, so it is zero. At the symmetric potential,
two copies of each old edge contribute half its old energy each,
and the diagonal edges contribute zero. Hence the defect is identical.

For the lower bound consider only realizations with uniform \(u\),
two nonempty cells, zero twists across cells, zero minimizing potential
near \(0\), and increments constant in each cell near \(0\).
The analytic functions \(\cosh(at)-1\) for distinct \(a>0\) are linearly
independent by their first even coefficients (a Vandermonde matrix).
Thus every vertex in a cell sees all \(s_\pm\) nonzero twist magnitudes
of that cell. Independence of the corresponding \(\sinh(at)\) makes
each magnitude balance separately, requiring at least two incident
edges per magnitude at each vertex. Each cell has at least \(2s_\pm+1\)
vertices. The difference of increments has at most \(s_++s_-\)
magnitudes. To begin at degree \(2m\) it needs at least \(m\)
magnitudes, by the same Vandermonde argument. T4 identifies this degree
when the defect has order \(4m\); thus \(n\geq2m+2\).
\end{proof}

\section{Minimum dimensions and constant-diagonal rigidity}
\begin{theorem}[T8]
Let \(N_m\) be the least dimension allowing a strictly positive leading
term of order \(4m\) under the initial assumptions. Then
\(N_1=3,\ N_2=3,\ N_3=4\). For \(m\geq4\),
\(4\leq N_m\leq2m+2\).
\end{theorem}
\begin{proof}
On a tree every field is a gradient, so dimension two is impossible.
A connected three-vertex support is a tree or a triangle. On a triangle
with nonzero circulation, if the first two tests from T5 are constant,
\(Q(x,y)\) has the same value on each pair of its three resistances.
For \(x\) the smallest, \(Q(x,s)\) is strictly increasing on \(s\geq x\);
hence the other two resistances are equal, say \(y\). Then
\(Q(x,y)=Q(y,y)\) gives \(x(x-y)=0\), forcing all equal. The first test
gives uniform \(u\), and T5 gives equality rather than order twelve
or higher. The three witnesses in Section 4 attain the lower bounds,
and T7 supplies the general upper bound.
\end{proof}

\begin{theorem}[T9]
On unicyclic support with constant diagonal, the only possibilities
are order four or identically zero, with the same nonzero-circulation
equality class as T5. If \(Z_m\) is the minimum for zero-diagonal
bipartite realizations of order \(4m\), then
\[
 Z_1=4,\qquad Z_2=5,\qquad 5\leq Z_3\leq8 .
\]
A five-vertex \(K_{2,3}\) witness has coefficient \(5/512\) at order eight.
\end{theorem}
\begin{proof}
A common diagonal can be subtracted without changing \(u,c,E\)
or the defect: both positive Perron eigenvalues decrease by exactly
that diagonal. Zero circulation and outside vertices are covered
by T3 and T5. On a zero-diagonal pure cycle with nonzero circulation
let \(J_0=\Phi/\sum_E c_e^{-1}\). If \(g=0\), the harmonic field
\(\alpha_e=J_0/c_e\) gives
\[
 J_0^2(c_{i-1}^{-1}+c_i^{-1})=2a u_i^2,\qquad
 \lambda_0 u_i^2=c_{i-1}+c_i .
\]
Cancel the positive sum to get a constant product
\(c_{i-1}c_i=\lambda_0J_0^2/(2a)\). Hence the conductances are
two-periodic, \(u\) is uniform and T5 gives equality. Otherwise T2
gives order four.

For the witnesses, use positive rational \(C_{ij}\) on edges, masses
\(\mu_i=\sum_j C_{ij}\), and \(U_i=\sqrt{\mu_i}\). Set
\(B_{ij}=C_{ij}/(U_iU_j)\) on edges, with zero diagonal, so \(BU=U\);
\(M=\sum_i\mu_i\) and \(u=U/\sqrt M\).
For a four-cycle use \(C=(1,1,1,2)\), with harmonic
\(h=(2/7,2/7,2/7,1/7)\). The masses are \((3,2,2,3)\).
The centered second-order increment is
\(f_2=(-2,3,3,-2)/245\). In these coordinates
\(\sum_j C_{ij}f_{2,j}/\mu_i=f_{2,i}/6\) and
\(\sum_i\mu_i f_{2,i}^2=60/245^2\).
The fourth-order coefficient is \((6/5)(60/245^2)/10=36/300125\).

For \(K_{2,3}\), orient its three two-edge paths from \(L\) to \(R\)
through \(M_1,M_2,M_3\). Give the six consecutive edges conductances
\((1,4,4,1,4,4)\) and twists \((2,1/2,1/2,2,-1,-1)\).
Path currents are \(2,2,-4\), so the twist is harmonic.
The masses are \((9,9,5,5,8)\), and \(k_2=1/2\) at every vertex.
The cubic potential coefficient \(\eta_3\) solves
\(L_C\eta_3=\operatorname{div}(C h^3)/6\), where \(L_C\) is the
positive weighted graph Laplacian and divergence sums outward fields.
Up to a constant it is \((5/24,-5/24,-3/8,3/8,0)\):
substitution on the three paths verifies the equation.
The fourth-order increments follow by summing
\(C_e[h_e\nabla_e\eta_3+h_e^4/24]/\mu_i\) at each vertex. They are
\((3/32,3/32,-1/32,-1/32,1/4)\), of weighted mean \(3/32\).
Thus \(f_4=(0,0,-1/8,-1/8,5/32)\). Each
\(\sum_j C_{ij}f_{4,j}=0\); consequently \(B(U f_4)=0\) and
\(R(U f_4)=U f_4\). The defect coefficient is
\(\sum_i\mu_i f_{4,i}^2/36=5/512\).

Connected bipartite supports on at most three vertices are trees,
and on four vertices the only non-tree possibility is a four-cycle.
The dichotomy excludes order eight or higher on that zero-diagonal
cycle. The witnesses attain \(Z_1=4\) and \(Z_2=5\), and lifting
the order-twelve four-cycle gives the upper bound \(Z_3\leq8\).
\end{proof}

\section{Scope and reproducibility}
The proof covers the precise revision-5 statements T1--T9 of
\texttt{C-e83ae1e1d8f8447c846f8208}. The general \(2m+2\) upper bound
does not prove global optimality, \(Z_3\) is not determined, and the
bipartite minima do not classify unrestricted zero-diagonal realizations.
All asymptotic remainders are for fixed \(B,h\).

The companion proof files include further exact certificates, including
a direct characteristic-polynomial check for a triangle with a leaf.
Five Python reproduction programs and the separate bounded search give
finite checks, not proofs of universal assertions. The exact clarified
proof snapshot passed two independent full paper reviews and a separate
bounded counterexample and boundary review
(fixed record \texttt{F-85c183ad590bceac57b760cdf542e402}).
The review and drafting process used AI assistance. This is a
versioned paper-track validation, with no Lean formalization or claim of
journal peer review. The article and original exposition are available
under CC BY 4.0; the supplied reproduction scripts are MIT-licensed.

\bibliographystyle{plain}
\bibliography{references}
\end{document}
